% Session 1 -- The TPS Report Generator (the reference session).
% Law 1 sources: os/milton/tpsrep_program.asm (command line from the PSP
% default FCB; TPSPARM.DAT keys read by getparam at line starts; LINES via
% atoi, 0 -> 1; the six required fields of reqtab; deleted records
% skipped; the cover / pagehead / detail / totals pieces; return codes
% 0/1/5/6/7/8 of finish); spec/tps_ritual.json (rows R1-R5, E6-E8,
% file_rules); harness/diff/tps/tpsmk.prg (the table, written by real dBASE
% III PLUS) and tpsrept_expected.txt (the report, real MS-DOS 6.22 golden,
% byte-equal to the flagship's by test-flair-tps-journey); Makefile
% flair-data rules (the three files in APPS on the shipped volume).
% GATES: test-tps-report test-flair-tps-journey test-pubs-tps test-pubs-screens
\Gates{test-tps-report, test-flair-tps-journey, test-pubs-tps, test-pubs-screens}
\chapter{The TPS Report Generator}

\SideHead{Introduction}This session describes the TPS Report Generator:
what it reads, what it writes, the order in which it requires the steps of
the operating procedure, and the report it produces.  The later sessions
teach each step as a numbered procedure.  Keep this session at hand as the
reference while you work.

\SideHead{Objectives}Upon completion of this session you will be able to:
\begin{enumerate}
\item Name the files the program reads and writes, and the directory that
  holds them.
\item State the form of the \cmd{TPSREP} command.
\item Name the four lines of the parameter record.
\item Name the six fields of a TPS table.
\item Name the five steps of the operating procedure in the order the
  program requires them.
\item Name the parts of the report and the return codes of the program.
\end{enumerate}

\section{What the Program Does}
\index{TPS Report Generator}
The TPS Report Generator reads a TPS table, a dBASE III table of TPS
reports on file, and writes the TPS report: a cover sheet, one or more
pages of detail lines, and the totals.  The report is displayed on the
screen as it is written and is kept in the file \rpt\ in the current
directory.

The program runs only from the InitechDOS prompt.  It is not started from
the FLAIR desktop: the operator leaves the desktop first (Session~2).

Before it writes anything, the program requires the steps of the operating
procedure, in this order:

\begin{center}\small
\begin{tabular}{@{}llp{48mm}@{}}
\rule{0pt}{10pt}\sffamily\bfseries Step & \sffamily\bfseries Procedure & \sffamily\bfseries What the program requires\\\hline
1 & \proc{02} & A parameter record, \file{TPSPARM.DAT}, in the current directory\\
2 & \proc{03} & The form aligned: the \key{Space Bar} pressed the required number of times, then \key{Enter}\\
3 & \proc{04} & The report number keyed twice, both times the same\\
4 & \proc{05} & The number of the memo about the new cover sheets\\
5 & \proc{06} & The attestation that the new cover sheet is attached\\\hline
\end{tabular}\end{center}

A step that is not satisfied is refused with a message, and the program
either asks again or ends.  Session~7 lists every refusal.

\section{Files Used}
\index{files used}
The program and its files are in the directory \file{\bs APPS} of the
shipped InitechOS data volume, drive A.

\begin{center}\small
\begin{tabular}{@{}lp{74mm}@{}}
\rule{0pt}{10pt}\sffamily\bfseries File & \sffamily\bfseries Contents\\\hline
\file{TPSREP.COM} & The TPS Report Generator\\
\file{TPSPARM.DAT} & The parameter record: report title, division, memo number and lines per page\\
\file{TPS.DBF} & The TPS table: the TPS reports on file\\
\rpt & The report, written by the program (cover sheet, detail pages, totals)\\
\file{TPSREPT.BAK} & The previous report, kept by the program\\\hline
\end{tabular}\end{center}

The program reads \file{TPSPARM.DAT} and the table from the current
directory, and writes \rpt\ and \file{TPSREPT.BAK} in the current directory.
For this reason the directory \file{A:\bs APPS} is made the current directory
before the program is run (\proc{01}).

\section{The TPSREP Command}
\index{TPSREP command}
\RegisterCommand{TPSREP}{external}{GATED}{TPSREP}%
\SideHead{Format}\SyntaxLine{TPSREP \ph{table}.DBF}

\SideHead{Remarks}\ph{table}\texttt{.DBF} is the name of the TPS table in
the current directory.  On the shipped volume the table is \file{TPS.DBF}.
The name is given as a file name only, without a drive or a path.  If no
table is named, or the named table cannot be opened, the program displays
\TpsLine{TPS-007}{1} and ends with return code 7.

The command is typed at the InitechDOS prompt and ended with \key{Enter}.
The program displays its banner, and then the first step of the operating
procedure (Session~4).

\SideHead{Example}With \file{A:\bs APPS} the current directory, the
operator types the command and the program displays its banner:

\OpsShot[fromtext={A:\Bsl APPS>TPSREP TPS.DBF}, totext={(C) 1989}]{tpsstart}

Typed without a table name, the command is refused:

\OpsShot[fromtext={A:\Bsl APPS>TPSREP}]{tps007}

\section{The Parameter Record}
\index{parameter record}\index{TPSPARM.DAT}
The parameter record \file{TPSPARM.DAT} is a text file of four lines.
Each line begins with its name, an equal sign, and the value; the value
runs to the end of the line.  The lines may be in any order.  The
parameter record on the shipped volume is:

\ReportLines[label=fig:parm, caption={The parameter record \file{TPSPARM.DAT} on the shipped volume.}]{../../../harness/diff/tps/tpsparm.dat}{1}{4}

\SideHead{TITLE=}The report title, printed centred on the cover sheet
under \texttt{T P S\ \ \ R E P O R T}.

\SideHead{DIVISION=}The division, printed on the cover sheet.

\SideHead{MEMO=}The number of the memo about the new cover sheets.  The
operator keys this number in \proc{05}; it is printed on the cover sheet
and in the totals.

\SideHead{LINES=}The number of detail lines on each page of the report.
A value of 0 is taken as 1.

A missing line is refused before anything else is done: the program
displays \TpsLine{TPS-002}{1}, with the name of the missing line, and
ends with return code 1.  The parameter record is maintained under
Operating Procedure OP-7.3, which is \proc{02} of this book.

\section{The TPS Table}
\index{TPS table}\index{TPS.DBF}
A TPS table is a dBASE III table with these six fields:

\begin{center}\small
\begin{tabular}{@{}llrp{46mm}@{}}
\rule{0pt}{10pt}\sffamily\bfseries Field & \sffamily\bfseries Type & \sffamily\bfseries Width & \sffamily\bfseries Contents\\\hline
\texttt{RPT\_NO} & Numeric & 6 & The number of the TPS report on file\\
\texttt{EMPLOYEE} & Character & 20 & The employee who filed it\\
\texttt{DEPT} & Character & 12 & The department\\
\texttt{PAGES} & Numeric & 3 & The number of pages filed\\
\texttt{COVER} & Logical & 1 & Whether the cover sheet was attached\\
\texttt{FILED} & Date & 8 & The date it was filed\\\hline
\end{tabular}\end{center}

A table without one of these fields is refused with \TpsLine{TPS-006}{1}
(return code 6).  A record marked for deletion is not listed.  The shipped
table \file{TPS.DBF} holds 14 records, one of them marked for deletion.

\section{The Report}
\index{report}\index{cover sheet}
The report has three parts.

\SideHead{Cover sheet}The cover sheet, \form{1} (Rev.\ 10/89), comes
first.  It carries the report number, the division, the name of the table,
the date the table was last updated, the tables on file in the current
directory, the memo number and the number of records in the table.

\ReportLines[label=fig:cover, caption={The cover sheet, as the program writes it to
\file{TPSREPT.TXT}.}]{../../../harness/diff/tps/tpsrept_expected.txt}{1}{17}

\SideHead{Detail pages}Each page begins with the report number and the
page number, then the column heads, then as many detail lines as the
\texttt{LINES=} line of the parameter record allows.  In \rpt\ each page
begins on a new sheet; on the screen a blank line separates the pages.

\ReportLines[caption={The first detail page (\texttt{LINES=6}).}]{../../../harness/diff/tps/tpsrept_expected.txt}{18}{26}

\SideHead{Totals}The totals give the reports listed, the pages filed and
the number of reports filed without a cover sheet.  When that number is
not zero, the totals end with the instruction to see the memo.

\OpsShot[fromtext={REPORTS LISTED}, totext={*** END OF TPS REPORT ***}]{tpsdone}

\section{Return Codes}
\index{return codes}\index{ERRORLEVEL}
The program ends with the line
\TpsMsg{END}
where \ph{n} is the return code.
The return code can be tested in a batch file with \cmd{IF ERRORLEVEL}.

\begin{center}\small
\begin{tabular}{@{}rp{88mm}@{}}
\rule{0pt}{10pt}\sffamily\bfseries Code & \sffamily\bfseries Meaning\\\hline
0 & The report was written\\
1 & The parameter record is missing or incomplete (TPS-001, TPS-002)\\
5 & The cover sheet was not attached; no report was written (TPS-005)\\
6 & The table is not a TPS table (TPS-006)\\
7 & No table was named, or it was not found (TPS-007)\\
8 & \rpt\ could not be created (TPS-008)\\\hline
\end{tabular}\end{center}

\section{Study Questions}
\begin{enumerate}
\item In which directory of the shipped volume are the TPS Report Generator
  and its files?
\item What are the four lines of the parameter record?
\item Which file holds the previous report?
\item In what order does the program require the five steps of the
  operating procedure?
\item What return code means that the report was written?
\end{enumerate}
